\section{הבעיה והזדמנות השוק \en{(Problem \& Market)}}

\subsection{הבעיה}
תצלומים מודפסים מתכלים: הצבעים דוהים, הנייר נסדק, כתמי לחות ושריטות פוגעים בפרטי הפנים. בו־זמנית, רוב המשפחות אינן יודעות מי מופיע בתצלומים הישנים. כאשר הדור המבוגר נפטר, השמות, המקומות והסיפורים נעלמים יחד איתו. הפתרונות הקיימים מתחלקים לשלושה סוגים, וכל אחד מהם חלקי.

\begin{table}[h]
\centering
\caption{השוואה בין חלופות קיימות לבין \en{Memoria}}
\label{tab:alternatives}
\small
\begin{tabularx}{\textwidth}{>{\bfseries}l X X X}
\toprule
\theadrow \thead{חלופה} & \thead{יתרון} & \thead{חסרון} & \thead{פער ש-\en{Memoria} סוגרת} \\
\midrule
סריקה עצמית בבית & זולה & איכות נמוכה, ללא שחזור & סריקה ארכיונית ושחזור \\
\en{Photoshop} ידני & תוצאה איכותית & יקר וארוך, כ-\num{2} שעות לתמונה & אוטומציה בסיוע \en{AI}, עם בקרת אדם \\
אפליקציות \en{AI} כלליות & מהירות & אין ארכיון, אין שירות, פרטיות לא ברורה & ארכיון מתועד והסכמה מפורשת \\
\bottomrule
\end{tabularx}
\end{table}

\subsection{פלח הלקוחות}
אנו מתמקדים בשלושה פלחים. הראשון הוא \emph{משפחות מורחבות} המתכננות אירוע (יום הולדת של סבתא, יום נישואין, שבעה) ורוצות אלבום ייחודי. השני הוא \emph{בני דור ביניים} בגילאי \num{40--60} המנקים בית של הורים. השלישי הוא \emph{ארגונים קטנים}: קיבוצים, מושבים, בתי כנסת ואגודות קהילתיות המחזיקים ארכיונים היסטוריים.

\subsection{אומדן שוק}
הערה מתודולוגית: אין בידינו נתון פומבי אמין על היקף שוק השחזור בישראל, ולכן אנו מציגים אומדן הנדסי ולא ציטוט. נניח \num{3.5} מיליון משקי בית בישראל, ושיעור קטן מהם, \num{0.5}\%, פונה לשירות מסוג זה בשנה. התוצאה היא כ-\num{17{,}500} לקוחות פוטנציאליים בשנה. מחיר הזמנה ממוצע של \num{\MAvgOrderPrice} ש״ח מציב את הפוטנציאל השנתי בסדר גודל של \num{13.7} מיליון ש״ח. שיעור החדירה שאנו מתכננים הוא כ-\num{\MOrdersEnd} הזמנות בחודש בסוף השנה השלישית, כלומר פחות מ-\num{15}\% מהפוטנציאל המשוער. נציין שהנחה זו נבדקת בניתוח הרגישות בפרק הפיננסי.

\subsection{יתרון תחרותי}
\begin{enumerate}[leftmargin=*]
  \item \textbf{שירות מלא} ולא כלי: סריקה, שחזור, תיעוד וגיבוי במקום אחד.
  \item \textbf{בקרת אדם} על כל תמונה לפני מסירה, כדי למנוע ״הזיות״ של מודל השחזור בפני אדם אמיתי.
  \item \textbf{פרטיות כברירת מחדל}: עיבוד מקומי, מחיקה לפי בקשה, והסכמה כתובה.
  \item \textbf{תמחור שקוף} ברמת הזמנה, ללא הפתעות.
\end{enumerate}

\subsection{דמויות לקוח}
\begin{figure}[h]
\centering
\begin{tikzpicture}[
  card/.style={draw=mteal,fill=mpaper,rounded corners=4pt,text width=0.30\textwidth,inner sep=7pt,font=\small,line width=0.8pt,align=center}]
  \node[card] (p1) {\textbf{\sffamily רונית, \num{58}}\\ סבתא לשמונה נכדים\\[2pt] רוצה אלבום ליום ההולדת של בעלה. מוכנה לשלם עבור איכות ושקט נפשי.};
  \node[card,right=4mm of p1] (p2) {\textbf{\sffamily דן, \num{47}}\\ מנקה את בית ההורים\\[2pt] מצא שלוש קופסאות תמונות ואינו יודע מי מופיע בהן. מחפש פתרון מקיף.};
  \node[card,right=4mm of p2] (p3) {\textbf{\sffamily ועד הקיבוץ}\\ ארכיון היסטורי\\[2pt] מעוניין בסריקה של \num{800} תמונות ובמטא־נתונים לשימוש ציבורי.};
  \foreach \p in {p1,p2,p3} \fill[mamber] ([yshift=-1.5pt]\p.north west) rectangle ([yshift=0pt]\p.north east);
\end{tikzpicture}
\caption{שלוש דמויות לקוח מייצגות, מבוססות על הפלחים שהוגדרו. הדמויות בדויות ונועדו לתכנון שירות.}
\label{fig:personas}
\end{figure}

כל דמות מובילה לחבילה אחרת: רונית ל-\en{Heritage} עם אלבום מודפס, דן ל-\en{Pro} בהיקף של כמה הזמנות, והקיבוץ להסכם ארגוני עם הנחת כמות. הדמויות משמשות אותנו גם כמבחן קבלה לתהליך: לכל אחת נכתב תסריט שימוש מקצה לקצה, וכל שלב בו נבדק מול ציפיות הלקוחה או הלקוח.
